% Industry CV of Vincenzo Di Florio (two pages at most).
% Build: cd cv && make   (writes ../docs/assets/pdf/CV_industry.pdf)
% Derived from cv-academic.tex and the April 2026 industry CV; same style (cv-style.sty).
\documentclass[10pt]{article}
\usepackage{cv-style}
\hypersetup{pdftitle={Vincenzo Di Florio — CV}, pdfauthor={Vincenzo Di Florio}}

\begin{document}

\cvheader{Computational scientist: numerical methods and scientific software in C++}

%==============================================================================
\section{Profile}
Physicist and mathematician (Ph.D.) specialising in numerical methods and scientific software for molecular simulation. Developed, with its co-authors, \textbf{NextGenPB}, a parallel finite element solver for the Poisson–Boltzmann equation built from scratch in C++, with applications to biomolecular free-energy calculations and drug-discovery workflows.

%==============================================================================
\section{Core skills}
\begin{itemize}
  \item \textbf{Programming}: C++ and Python (daily; numerical methods, scientific code, data analysis, Python interfaces to C/C++ codes), MATLAB / Octave (daily), Bash (scripting), Fortran (reading).
  \item \textbf{Parallel computing}: MPI / OpenMP (working knowledge, used to parallelise NextGenPB); Principal Investigator of a cloud computing grant (PRACE-FENIX, GAIA Cloud).
  \item \textbf{Numerical methods}: finite elements, adaptive and local mesh refinement on octree meshes, linear and nonlinear PDE solvers, sparse identification of dynamical systems (SINDy).
  \item \textbf{Machine learning}: physics-informed and physics-guided neural networks (PINNs, PGNNs) and data-driven resolution enhancement, in Python with PyTorch.
  \item \textbf{Domain}: biomolecular electrostatics, protein–protein binding, p\textit{K}\textsubscript{a} prediction, statistical mechanics.
  \item \textbf{Languages}: Italian (native), English (fluent).
\end{itemize}

%==============================================================================
\section{Experience}
\entry{Post-doctoral Research Fellow, MOX – Modeling and Scientific Computing}{Jan 2025 – present}{Department of Mathematics, Politecnico di Milano, Milan, Italy}
\begin{itemize}
  \item Develop numerical methods and adaptive mesh refinement strategies for the Poisson–Boltzmann equation, targeting super-resolution accuracy in molecular electrostatics.
  \item Study protein–protein binding through electrostatic complementarity at biomolecular interfaces (poster at EuroQSAR 2026).
  \item Design physics-guided machine learning models that accelerate electrostatic potential calculations, including data-driven resolution enhancement of NextGenPB potentials (talk at GIMC SIMAI Young 2026).
  \item Teaching assistant in Advanced Methods for Scientific Computing (C++ for scientific computing); tutor in Mathematical Physics.
\end{itemize}

\entry{Research collaboration, CONCEPT Lab, Istituto Italiano di Tecnologia}{Nov 2024 – Dec 2024}{Genoa, Italy}
\begin{itemize}
  \item Evaluated and selected the numerical model for a nonlinear Poisson–Boltzmann solver; carried out the feasibility analysis and benchmarked candidate discretisation strategies.
\end{itemize}

\entry{Internship, City College of New York (CUNY)}{May 2024 – Aug 2024}{New York, USA}
\begin{itemize}
  \item Integrated NextGenPB as the electrostatic solver of the MCCE multi-conformation continuum electrostatics workflow, used for biomolecular free-energy calculations.
  \item Built and validated a benchmark pipeline for p\textit{K}\textsubscript{a} prediction, comparing physics-based results with experimental reference data for a set of proteins.
\end{itemize}

\entry{Ph.D. research, Politecnico di Torino and Istituto Italiano di Tecnologia}{2021 – 2025}{Turin and Genoa, Italy}
\begin{itemize}
  \item NextGenPB: finite elements, analytical super-resolution, local mesh (de)refinement, parallel execution. Published in Computer Physics Communications (2025, first author) and released as open source (\href{https://github.com/vdiflorio/NextGenPB}{github.com/vdiflorio/NextGenPB}).
  \item Teaching assistant in Calculus 1 and Physics 1; tutor in Methods for Environmental Engineering.
\end{itemize}

\entry{Internship, SISSA}{Sep 2018 – Dec 2018}{Trieste, Italy}
\begin{itemize}
  \item Applied sparse regression (SINDy) to identify the governing equations of dynamical systems from time-series data.
\end{itemize}

%==============================================================================
\section{Education}
\entry{Ph.D. in Pure and Applied Mathematics}{2021 – 2025}{Politecnico di Torino and Istituto Italiano di Tecnologia}
\entry{Master's degree in Physics of Complex Systems (110/110 cum laude)}{2019 – 2021}{Università degli Studi di Torino}
\entry{Bachelor's degree in Physics (110/110 cum laude)}{2016 – 2019}{Università degli Studi di Trieste}

%==============================================================================
\section{Publications and talks}
Six peer-reviewed journal articles (2024–2026); the full list is on \href{https://vdiflorio.github.io}{vdiflorio.github.io}. Selected:
\begin{itemize}
  \item \textbf{Di Florio, V.} et al. (2025). \emph{NextGenPB: an analytically-enabled super resolution tool for solving the Poisson–Boltzmann equation featuring local (de)refinement.} Computer Physics Communications. \doi{10.1016/j.cpc.2025.109816}
  \item Mazzeo, M. D., \textbf{Di Florio, V.}, et al. (2026). \emph{Patch-based ray tracing in NanoShaper boosts molecular surface computation.} Journal of Chemical Information and Modeling. \doi{10.1021/acs.jcim.5c02287}
  \item Alharbi, S., …, \textbf{Di Florio, V.}, et al. (2026). \emph{SCRIB inhibits PP1 phosphatase activity without blocking or structurally altering its active site.} Journal of Biological Chemistry. \doi{10.1016/j.jbc.2026.113599}
\end{itemize}
Talks at SIAM PP26 (Berlin, 2026), GIMC SIMAI Young (Pisa, 2026) and ICIAM (Tokyo, 2023); posters at EuroQSAR 2026, Protein Electrostatics 2025, ACS Fall 2024 and Protein Electrostatics 2023.

%==============================================================================
\section{Authorization to personal data processing}
{\small\color{muted}
\begin{otherlanguage}{italian}
Autorizzo il trattamento dei dati personali contenuti nel mio curriculum vitae ai sensi del GDPR 2016/679 e del D.~Lgs.~196/2003.\par
\end{otherlanguage}}

\end{document}
